\chapter{Stage Appendix: Part I -- Parent Theory, Geometry Lift, and Linearized PDE Backbone}
\label{app:stage-part-1}

\section*{Stages 001--006; anchors MTDC-T1--T4}
\addcontentsline{toc}{section}{Stages 001--006; anchors MTDC-T1--T4}

This appendix is the detailed verification ledger for the first theorem block of the moving-throat derivation.  The corresponding main-text chapter, \cref{chap:part01_parent_geometry}, explains the role of the block in the larger paper.  The job of this appendix is narrower and more mechanical: it records, stage by stage, the definitions, reductions, algebraic transforms, branch tests, and audit checks that turn the exact parent model into the first usable wall/support/outgoing response bundle.

The six stages here should be read as a single construction.  Stage~001 promotes the old collective geometry variables to a distributed moving-surface field and introduces the scalar, grouped real \(P_2\), and higher-harmonic lanes.  Stage~002 verifies that this distributed wall theory reduces back to the old \((a,L)\) closure in the lowest axisymmetric sector.  Stage~003 couples the wall to stable BdG matter support modes and derives the first conservative Schur-complement kernels.  Stage~004 now treats the electromagnetic sector projection-first: it derives the exact projected localized Maxwell law, identifies the matched zero-mode reduction, transports mouth-local projected corrections into the grouped \(Z_n,N_n\) slots, imports the parent-throat action packet, and retains the previous Maxwell/mixed one-port calculation as the reduced normal form for outgoing transfer.  Stage~005 translates that outgoing bridge into the normalized grouped-response language used by the retarded quadrupole program.  Stage~006 writes the full grouped \(20/21/22\) bundle, its weighted projectors, its microscopic coefficient ledger, and the isotropic normalization test.

\paragraph{Source layout.}
Stages~001--006 are now sourced from the canonical files \StageFile{stages/stage_001.tex} through \StageFile{stages/stage_006.tex}.  The raw Markdown notes and other provenance artifacts are tracked in a repo-local provenance index, so they do not have to appear in the reader-facing PDF.  The stage files themselves carry the executable SymPy, Mathematica, and numerical-stress references when they exist.

\begin{longtable}{@{}L{0.10\textwidth}L{0.27\textwidth}L{0.22\textwidth}L{0.34\textwidth}@{}}
\caption{Part I stage-audit map.}\label{tab:app-part01-stage-map}\\
\toprule
Stage & Source stage & Primary status & Verification output \\
\midrule
\endfirsthead
\toprule
Stage & Source stage & Primary status & Verification output \\
\midrule
\endhead
001 & Geometry lift and linearized PDE skeleton & \StatusEffective{} / \StatusExactClosure{} & Distributed shape field \(R(\Omega,w,t)\), promoted confinement coupling, modal split, and first linearized GNLS/Maxwell/wall system. \\
\addlinespace
002 & Breathing reduction back to \((a,L)\) & \StatusExactClosure{} & Two-mode axisymmetric reduction, effective \((\delta a,\delta L)\) mass/stiffness matrices, and uncoupled grouped \(P_2\) degeneracy. \\
\addlinespace
003 & Minimal BdG--wall coupling & \StatusReduced{} / \StatusExactClosure{} & Stable-mode Schur complement, low-frequency wall moments, pole shift formula, and conservative grouped-isotropy diagnostic. \\
\addlinespace
004 & Projection-first Maxwell/mixed response & \StatusExact{} / \StatusReduced{} / \StatusExactClosure{} & Exact projected localized Maxwell law, matched reduction channel, projected \(Z_n,N_n\) bundle transport, reduced one-port transfer factor, compact \(l=2\) fingerprint, and scalar-compatibility criterion. \\
\addlinespace
005 & Grouped normalization bridge & \StatusExactClosure{} / \StatusOpen{} & Conversion from operator moments to normalized response moments and invariant outgoing-normalization product. \\
\addlinespace
006 & Full grouped bundle and projectors & \StatusExactClosure{} / \StatusOpen{} & Weighted grouped projectors, full \(D,N,P\) coefficient ledger, isotropic branch test, and first-order anisotropy transport laws. \\
\bottomrule
\end{longtable}

\paragraph{Audit philosophy.}
The derivations below preserve the status distinction used throughout the ledger.  Algebraic eliminations, projector identities, low-frequency expansions, and grouped inverse maps are exact once the declared reduced variables are adopted.  The projection-first Maxwell law is exact after the observation kernel is declared; the matched zero-mode and one-port Maxwell/mixed normal forms are controlled reductions.  The adoption of a distributed wall closure, a stable BdG normal-mode reduction, and a compact outgoing port are reduced or effective moves.  The final equality of the outgoing-normalization product with the universal quadrupole target is not proved in this appendix; it is formulated here as the branch test that later stages attempt to realize.

\paragraph{Notation warning.}
The grouped symbols \(20,21,22\) are grouped real \(P_2\) lanes, not spacetime indices.  The \(21\) and \(22\) labels each abbreviate a cosine/sine real pair when the full five-dimensional real STF representation is being discussed.  The mixed gauge variables \(A_w,J^w,F_{\mu w},E_w,C_a\) are retained as microscopic channels; they are suppressed only in the strict far-field brane reduction.

\begingroup
\let\clearpage\relax
\clearpage
\section[Stage 001]{Stage 001: Geometry Lift and Linearized PDE Skeleton}
\label{stage:001}

\stagefield{Part / anchor}{Part I; primary anchors \resultanchor{MTDC-T1}, \resultanchor{MTDC-T2}, and \resultanchor{MTDC-T4}.}
\claimstatus{The parent-field import is \StatusExact{} inside the declared parent model.  The shape-field lift and minimal quadratic wall action are \StatusEffective{}.  The modal split and the linearized equations are \StatusExactClosure{} once that lift is adopted.}

\stagefield{Purpose}{Stage~001 replaces the old finite-dimensional geometry variables \(a(t)\) and \(L(t)\) by a distributed throat geometry variable, while preserving those variables as collective moments.  The stage also introduces the first linearized matter/gauge/geometry system and the harmonic decomposition in which the scalar lane, grouped real \(P_2\) lane, geometry lane, and higher modes become explicit.}

\stagefield{Inputs}{The stage imports the exact \((4+1)\)-dimensional parent theory: gauged GNLS matter, localized Maxwell field, corrected charge ontology, and confinement potential.  It also imports the lower-order requirement that the moving-throat PDE expose the already-visible \(P_0\oplus P_2\) support sector, the separate geometry-completion channel, and the mixed Maxwell channels retained outside the strict zero-mode brane limit.}

\stagefield{Verification}{SymPy audit: \StageFile{scripts/moving_throat_pde_stage001_geometry_lift_sympy_audit.py}.  Mathematica audit: \StageFile{mathematica/moving_throat_pde_stage001_geometry_lift_mathematica_audit.wl}.  The audits check the monopole/\(P_2\) harmonic bookkeeping, the confinement chain-rule sign, and the densitized-versus-weighted wall-action convention used by the ledger.}

\subsection*{Derivation}

\paragraph{1. Shape-field lift.}
Let
\begin{equation}
\label{eq:app-stage001-sigma}
\Sigma(\mathbf X,t)=r-R(\Omega,w,t),
\qquad
r=\sqrt{x^2+y^2+z^2},
\qquad
\Omega=\mathbf x/r\in S^2.
\end{equation}
The finite throat surface is the level set
\begin{equation}
\label{eq:app-stage001-surface}
\Sigma(\mathbf X,t)=0.
\end{equation}
The working sign convention is
\begin{equation}
\Sigma>0\quad\text{outside the throat},
\qquad
\Sigma<0\quad\text{inside the support/interior region}.
\end{equation}
This choice keeps the level-set coupling to the bulk fields while making the mouth-sphere harmonics explicit.  It is the minimum lift needed to avoid treating the grouped real \(P_2\) coefficients as purely symbolic residuals.

\paragraph{2. Reference stationary branch.}
Take a stationary reference surface
\begin{equation}
\label{eq:app-stage001-reference}
\Sigma_0(\mathbf X)=r-R_0(w)=0,
\end{equation}
with mouth at \(w=0\), finite depth \(0\le w\le L_0\), mouth radius
\begin{equation}
\label{eq:app-stage001-a0}
a_0=R_0(0),
\end{equation}
and a bottom closure such as \(R_0(L_0)=0\) or an equivalent regularity condition.  The old collective variables are recovered as lowest moments:
\begin{equation}
\label{eq:app-stage001-a-moment}
a(t)=\frac{1}{4\pi}\int_{S^2}R(\Omega,0,t)\,d\Omega,
\end{equation}
and, if \(W_b(\Omega,t)\) is defined implicitly by
\begin{equation}
R(\Omega,W_b(\Omega,t),t)=0,
\end{equation}
then
\begin{equation}
\label{eq:app-stage001-l-moment}
L(t)=\frac{1}{4\pi}\int_{S^2}W_b(\Omega,t)\,d\Omega.
\end{equation}
Thus the old closure survives as a moment projection of the distributed wall field.

\paragraph{3. Promoted confinement coupling.}
The old confinement dependence \(V_{\rm conf}(\mathbf X;a,L)\) is promoted to a moving-surface dependence
\begin{equation}
\label{eq:app-stage001-promoted-conf}
V_{\rm conf}(\mathbf X;\Sigma)=V_{\rm wall}\!\left(\frac{\Sigma(\mathbf X,t)}{\ell_c}\right),
\end{equation}
where \(\ell_c\) is the wall thickness.  Around the reference branch,
\begin{equation}
V_0(\mathbf X)=V_{\rm wall}\!\left(\frac{\Sigma_0(\mathbf X)}{\ell_c}\right).
\end{equation}
The first variation is
\begin{equation}
\label{eq:app-stage001-delta-v-general}
\delta V_{\rm conf}
=\frac{V_{\rm wall}'(\Sigma_0/\ell_c)}{\ell_c}\,\delta\Sigma.
\end{equation}
Since \(\Sigma=r-R\), a wall perturbation \(R=R_0+\eta\) gives
\begin{equation}
\label{eq:app-stage001-delta-sigma}
\delta\Sigma=-\eta(\Omega,w,t),
\end{equation}
and therefore
\begin{equation}
\label{eq:app-stage001-delta-v}
\boxed{
\delta V_{\rm conf}
=-\frac{V_{\rm wall}'(\Sigma_0/\ell_c)}{\ell_c}\,\eta(\Omega,w,t)}.
\end{equation}
This is the first direct wall-to-BdG coupling used later.

\paragraph{4. Harmonic decomposition.}
Write
\begin{equation}
\label{eq:app-stage001-wall-perturbation}
R(\Omega,w,t)=R_0(w)+\eta(\Omega,w,t).
\end{equation}
Expand \(\eta\) in real spherical harmonics:
\begin{equation}
\label{eq:app-stage001-harmonic-expansion}
\eta(\Omega,w,t)
=\eta_0(w,t)Y_{00}(\Omega)
+\sum_{m\in P_2({\rm real})}q_{2m}(w,t)Y^{\rm real}_{2m}(\Omega)
+\eta_{\ge3}(\Omega,w,t).
\end{equation}
The real \(P_2\) set is
\begin{equation}
\label{eq:app-stage001-p2-set}
P_2({\rm real})=\{20,21c,21s,22c,22s\}.
\end{equation}
With the normalized monopole
\begin{equation}
Y_{00}=\frac{1}{2\sqrt\pi},
\qquad
\int_{S^2}Y_{00}^2\,d\Omega=1,
\end{equation}
we have
\begin{equation}
\frac{1}{4\pi}\int_{S^2}Y_{00}\,d\Omega=\frac{1}{2\sqrt\pi}.
\end{equation}
Therefore the normalized monopole coefficient and the physical mouth-average shift are related by
\begin{equation}
\label{eq:app-stage001-q00-da}
q_{00}(0,t)=2\sqrt\pi\,\delta a(t).
\end{equation}
A useful axisymmetric split is
\begin{equation}
\label{eq:app-stage001-axisymmetric-split}
\eta_0(w,t)=\alpha_a(w)\delta a(t)+\alpha_L(w)\delta L(t)+g(w,t),
\end{equation}
where \(g\) is the residual axisymmetric geometry lane orthogonal to the \((a,L)\) collective moments.

\paragraph{5. Minimal distributed wall action.}
The working passive quadratic wall action is
\begin{equation}
\label{eq:app-stage001-wall-action}
S_{\eta}^{(2)}
=\frac12\int dt\,dw\,d\Omega\,\sqrt{\gamma_0}
\left[
\mu_\eta(w)(\partial_t\eta)^2
-T_w(w)(\partial_w\eta)^2
-T_\Omega(w)\eta(-\Delta_{S^2})\eta
-K_\eta(w)\eta^2
\right].
\end{equation}
The coefficients \(\mu_\eta,T_w,T_\Omega,K_\eta\) are branch constitutive functions.  From this point onward the ledger adopts the densitized one-dimensional convention: the background surface weight \(\sqrt{\gamma_0}\) is absorbed into the effective axial coefficients and mode profiles, so the reduced formulas are written with respect to \(dw\).  The weighted surface form is still legitimate, but then the Euler--Lagrange operator carries the extra first-derivative term coming from \(\partial_w\sqrt{\gamma_0}\); the SymPy and Mathematica audits check both forms explicitly.

Projecting onto a harmonic \(Y_{lm}\), using
\begin{equation}
-\Delta_{S^2}Y_{lm}=l(l+1)Y_{lm},
\end{equation}
gives the modal wall operator
\begin{equation}
\label{eq:app-stage001-modal-wall-pde}
\boxed{
\mu_\eta\,\partial_t^2q_{lm}
-\partial_w\!\bigl(T_w\partial_wq_{lm}\bigr)
+\bigl[K_\eta+l(l+1)T_\Omega\bigr]q_{lm}
=S_{lm}^{(\psi,A)}+f_{lm}^{\rm ext}}.
\end{equation}
The scalar lane \(l=0\) and the grouped real quadrupole lane \(l=2\) are therefore separated before any additional closure.

\paragraph{6. Linearized matter and gauge sectors.}
Choose stationary background fields
\begin{equation}
\psi=e^{-i\mu_0t/\hbar}\psi_0(\mathbf X),
\qquad
A_M=A_{M0}(\mathbf X),
\qquad
R=R_0(w).
\end{equation}
Let
\begin{equation}
\psi=e^{-i\mu_0t/\hbar}\bigl(\psi_0+\delta\psi\bigr),
\qquad
A_M=A_{M0}+\delta A_M.
\end{equation}
The matter perturbation satisfies a BdG-type system
\begin{equation}
\label{eq:app-stage001-bdg-schematic}
i\hbar\partial_t
\begin{pmatrix}\delta\psi\\ \delta\psi^*\end{pmatrix}
=\mathcal L_{\rm BdG}
\begin{pmatrix}\delta\psi\\ \delta\psi^*\end{pmatrix}
+\mathcal C_A[\delta A_M]
+\mathcal C_\eta[\eta],
\end{equation}
where \(\mathcal C_\eta\) contains the coupling \eqref{eq:app-stage001-delta-v}.  The localized Maxwell equation linearizes to
\begin{equation}
\label{eq:app-stage001-linear-maxwell}
\boxed{
\partial_M\bigl(Z(w)\delta F^{MN}\bigr)
+\frac1\xi\partial^N(\partial\!\cdot\!\delta A)
=\mu_0\delta J^N}.
\end{equation}
At this stage the mixed variables
\begin{equation}
\delta A_w,
\qquad
\delta J^w,
\qquad
\delta F_{\mu w}
\end{equation}
are retained.  Their later suppression in the far-field brane limit is a reduction, not an ontology deletion.

\paragraph{7. Mouth/worldtube response operator.}
At the mouth \(w=0\), choose the real port basis
\begin{equation}
P_{00}=Y_{00},
\qquad
P_{20},P_{21c},P_{21s},P_{22c},P_{22s}.
\end{equation}
Let \(u_B(\omega)\) denote projected effort variables and \(j_A(\omega)\) the corresponding generalized fluxes or tractions.  The reduced response operator is written
\begin{equation}
\label{eq:app-stage001-response-operator}
j_A(\omega)=\sum_B Z_{{\rm eff},AB}(\omega)u_B(\omega).
\end{equation}
On an isotropic reference throat, this operator block-diagonalizes into scalar/geometry, grouped real \(P_2\), and higher-\(l\) sectors.

\stagefield{Output}{The stage outputs the shape-field closure \eqref{eq:app-stage001-sigma}, the confinement variation \eqref{eq:app-stage001-delta-v}, the harmonic split \eqref{eq:app-stage001-harmonic-expansion}, the modal wall PDE \eqref{eq:app-stage001-modal-wall-pde}, and the response-operator target \eqref{eq:app-stage001-response-operator}.  These are the definitions without which later grouped response coefficients have no PDE-side interpretation.}

\stagefield{Checks}{\begin{verificationchecklist}
\item The moment check follows from the normalized monopole integral and gives \(q_{00}=2\sqrt\pi\,\delta a\).
\item The sign check for \(\delta V_{\rm conf}\) follows from \(\delta\Sigma=-\eta\).
\item The harmonic selection check follows from orthogonality on \(S^2\), so the isotropic linear problem has no \(l=0\leftrightarrow l=2\) mixing.
\item The ontology check keeps \(A_w,J^w,F_{\mu w}\) alive before any far-field brane reduction.
\end{verificationchecklist}}

\stagefield{Downstream use}{Stage~002 uses the same wall action to recover \((a,L)\).  Stage~003 uses \eqref{eq:app-stage001-delta-v} as the wall--BdG coupling.  Stages~004--021 use the retained mixed variables as the port-active gauge sector and reduce them to the one-port normal form.  Stages~022--023 use the grouped real \(P_2\) response operator to formulate the outgoing-normalization product.}

\clearpage
\section[Stage 002]{Stage 002: Breathing Reduction Back to the Old \((a,L)\) Closure}
\label{stage:002}

\stagefield{Part / anchor}{Part I; primary anchors \resultanchor{MTDC-T1}, \resultanchor{MTDC-T2}, and \resultanchor{MTDC-T4}.}
\claimstatus{\StatusExactClosure{} inside the Stage~001 quadratic wall closure.  Damping and open-system completion are not derived here.}

\stagefield{Purpose}{Stage~002 verifies that the distributed wall lift is a true lift of the old geometry sector rather than a replacement unrelated to it.  The axisymmetric lowest-mode truncation reproduces a two-coordinate \((\delta a,\delta L)\) matrix system, while the grouped real \(P_2\) modes appear as the next harmonic family of the same wall PDE.}

\stagefield{Inputs}{The inputs are the normalized monopole harmonic, the Stage~001 wall action, the axisymmetric split of \(\eta\), and the grouped real \(P_2\) harmonic eigenvalue \(-\Delta_{S^2}Y_{2m}=6Y_{2m}\).}

\stagefield{Verification}{SymPy audit: \StageFile{scripts/moving_throat_pde_stage002_breathing_reduction_sympy_audit.py}.  Mathematica audit: \StageFile{mathematica/moving_throat_pde_stage002_breathing_reduction_mathematica_audit.wl}.  The audits check the \(Y_{00}\) normalization bridge, the derived \(4\pi\) overlap factor, the conservative \((\delta a,\delta L)\) matrix reduction, and the uncoupled grouped real \(P_2\) degeneracy.}

\subsection*{Derivation}

\paragraph{1. Axisymmetric two-mode ansatz.}
Using
\begin{equation}
Y_{00}=\frac{1}{2\sqrt\pi},
\qquad
\int_{S^2}Y_{00}^2\,d\Omega=1,
\qquad
\frac{1}{4\pi}\int_{S^2}Y_{00}\,d\Omega=\frac{1}{2\sqrt\pi},
\end{equation}
the axisymmetric perturbation is truncated as
\begin{equation}
\label{eq:app-stage002-two-mode-ansatz}
\eta_0(w,t)=2\sqrt\pi\left[\alpha_a(w)\delta a(t)+\alpha_L(w)\delta L(t)\right].
\end{equation}
Set
\begin{equation}
Q^A=(\delta a,\delta L),
\qquad
A\in\{a,L\}.
\end{equation}

\paragraph{2. Reduction of the quadratic wall action.}
On the axisymmetric branch, the angular-stiffness term vanishes because \(l=0\).  In the densitized one-dimensional convention inherited from Stage~001, the background surface weight has already been absorbed into the effective axial coefficients and profiles, so the reduction is written with respect to \(dw\) rather than \(\sqrt{\gamma_0}\,dw\).  The action therefore reduces to a Lagrangian
\begin{equation}
\label{eq:app-stage002-lred}
L_{\rm red}^{(0)}
=\frac12 M_{AB}\dot Q^A\dot Q^B
-\frac12 K_{AB}Q^AQ^B,
\end{equation}
where
\begin{equation}
\label{eq:app-stage002-mass-matrix}
\boxed{
M_{AB}=4\pi\int dw\,\mu_\eta(w)\alpha_A(w)\alpha_B(w)},
\end{equation}
and
\begin{equation}
\label{eq:app-stage002-stiffness-matrix}
\boxed{
K_{AB}=4\pi\int dw\left[
T_w(w)\alpha_A'(w)\alpha_B'(w)+K_0(w)\alpha_A(w)\alpha_B(w)
\right]}.
\end{equation}
Here \(K_0(w)=K_\eta(w)\) on the scalar branch.

\paragraph{3. Euler--Lagrange equations.}
Varying \eqref{eq:app-stage002-lred} gives
\begin{equation}
\label{eq:app-stage002-el}
\boxed{
M_{AB}\ddot Q^B+K_{AB}Q^B=0}.
\end{equation}
This is the conservative linear version of the old two-coordinate geometry closure.  The old schematic form
\begin{equation}
M_{AB}\ddot Q^B+\Gamma_{AB}\dot Q^B=-\frac{\partial H_{\rm tot}}{\partial Q^A}
\end{equation}
is recovered only after adding open-system or effective damping data.  Stage~002 derives the conservative part, not the damping matrix.

\paragraph{4. Grouped real \(P_2\) one-mode reduction.}
For one real quadrupole component take
\begin{equation}
\label{eq:app-stage002-p2-ansatz}
\eta_{2m}(\Omega,w,t)=\beta_2(w)q_{2m}(t)Y_{2m}^{\rm real}(\Omega),
\qquad
-\Delta_{S^2}Y_{2m}^{\rm real}=6Y_{2m}^{\rm real}.
\end{equation}
Inserting into the same wall action gives
\begin{equation}
\label{eq:app-stage002-l2m}
L_{2m}=\frac12 M_2\dot q_{2m}^{\,2}-\frac12K_2q_{2m}^{\,2},
\end{equation}
with
\begin{equation}
\label{eq:app-stage002-m2}
\boxed{
M_2=\int dw\,\mu_\eta(w)\beta_2(w)^2},
\end{equation}
and
\begin{equation}
\label{eq:app-stage002-k2}
\boxed{
K_2=\int dw\left[
T_w(w)\beta_2'(w)^2+\bigl(K_\eta(w)+6T_\Omega(w)\bigr)\beta_2(w)^2
\right]}.
\end{equation}
Thus, before matter/gauge coupling or explicit anisotropy,
\begin{equation}
\label{eq:app-stage002-p2-eom}
\boxed{
M_2\ddot q_{2m}+K_2q_{2m}=0
}
\end{equation}
for every real \(P_2\) component.  This is the first degeneracy theorem for the grouped real quadrupole lane.

\stagefield{Output}{The stage outputs the \((\delta a,\delta L)\) matrix system \eqref{eq:app-stage002-el}, the effective matrices \eqref{eq:app-stage002-mass-matrix}--\eqref{eq:app-stage002-stiffness-matrix}, and the uncoupled grouped \(P_2\) mass/stiffness pair \eqref{eq:app-stage002-m2}--\eqref{eq:app-stage002-k2}.}

\stagefield{Checks}{\begin{verificationchecklist}
\item The factor \(4\pi\) in \eqref{eq:app-stage002-mass-matrix} and \eqref{eq:app-stage002-stiffness-matrix} follows from the \(2\sqrt\pi\) normalization in \eqref{eq:app-stage002-two-mode-ansatz}.
\item The \(l=2\) angular stiffness shift is \(l(l+1)T_\Omega=6T_\Omega\).
\item The grouped \(P_2\) degeneracy follows because \(K_2\) and \(M_2\) do not depend on the real component label on an isotropic reference throat.
\item The reduction preserves the old \((a,L)\) closure but does not derive the old damping matrix.
\end{verificationchecklist}}

\stagefield{Downstream use}{Stage~003 couples the Stage~002 wall modes to stable BdG support modes.  Stage~023 uses the same wall baseline as the \(K_A,M_A\) entries of the full grouped bundle.  Later geometry-contamination stages use the distinction between collective scalar geometry and grouped \(P_2\) wall modes.}

\clearpage
\section[Stage 003]{Stage 003: Minimal BdG--Wall Coupling and First Pole Shifts}
\label{stage:003}

\stagefield{Part / anchor}{Part I; primary anchors \resultanchor{MTDC-T2} and \resultanchor{MTDC-T4}.}
\claimstatus{The Schur complement and pole formulas are \StatusExactClosure{} inside the declared stable-mode reduction.  Passing from the full BdG PDE to stable normal coordinates is \StatusReduced{}.}

\stagefield{Purpose}{Stage~003 turns on the first conservative matter-support dynamics.  It reduces the linearized BdG sector to stable normal coordinates, couples them to the wall through the confinement variation from Stage~001, integrates them out exactly, and obtains the first microscopic origin of wall pole shifts and grouped real \(P_2\) conservative response coefficients.}

\stagefield{Inputs}{The inputs are the Stage~001 coupling \(\delta V_{\rm conf}=-(V_{\rm wall}'/\ell_c)\eta\), the Stage~002 wall coordinates, the assumption of stable BdG normal modes with positive frequencies, and harmonic selection on an isotropic reference throat.}

\stagefield{Verification}{SymPy audit: \StageFile{scripts/moving_throat_pde_stage003_bdg_sympy_audit.py}; transcript: \StageFile{scripts/output/moving_throat_pde_stage003_bdg_sympy_audit.txt}.  Mathematica audit: \StageFile{mathematica/moving_throat_pde_stage003_bdg_mathematica_audit.wl}; transcript: \StageFile{mathematica/output/moving_throat_pde_stage003_bdg_mathematica_audit.txt}.  Shared numerical stress test with Stage~004: \StageFile{scripts/numerical/stage003_004_foundational_stress.py}, \StageFile{scripts/numerical/output/stage003_004_foundational_stress.txt}, \StageFile{mathematica/numerical/stage003_004_foundational_stress.wl}, and \StageFile{mathematica/numerical/output/stage003_004_foundational_stress.txt}.}

\subsection*{Derivation}

\paragraph{1. Stable scalar-sector reduction.}
Let
\begin{equation}
Q^A=(\delta a,\delta L),
\qquad
A\in\{a,L\},
\end{equation}
and let \(X_\alpha\) be stable scalar BdG normal coordinates with frequencies \(\varpi_{0\alpha}>0\).  The reduced scalar Lagrangian is
\begin{equation}
\label{eq:app-stage003-scalar-lagrangian}
L_{\rm red}^{(0)}=
\frac12\dot Q^TM_0\dot Q-\frac12Q^TK_0Q
+\frac12\sum_\alpha\bigl(\dot X_\alpha^2-\varpi_{0\alpha}^2X_\alpha^2\bigr)
+\sum_{A,\alpha}C_{A\alpha}Q^AX_\alpha.
\end{equation}
The matrix \(C\) consists of overlap integrals generated by \eqref{eq:app-stage001-delta-v} and the BdG mode profiles.

In frequency space,
\begin{equation}
(\varpi_{0\alpha}^2-\omega^2)X_\alpha=C_{A\alpha}Q^A,
\qquad
X_\alpha=(\varpi_{0\alpha}^2-\omega^2)^{-1}C_{A\alpha}Q^A.
\end{equation}
Substitution gives the exact scalar effective kernel
\begin{equation}
\label{eq:app-stage003-scalar-kernel}
\boxed{
D_0^{\rm eff}(\omega)
=K_0-\omega^2M_0-C(\Omega_0^2-
\omega^2I)^{-1}C^T},
\end{equation}
where
\begin{equation}
\Omega_0^2=\operatorname{diag}(\varpi_{0\alpha}^2).
\end{equation}
The low-frequency expansion is
\begin{equation}
\label{eq:app-stage003-scalar-expansion}
D_0^{\rm eff}(\omega)
=K_0^{\rm eff}-\omega^2M_0^{\rm eff}-\omega^4N_0^{\rm eff}+O(\omega^6),
\end{equation}
with
\begin{equation}
\label{eq:app-stage003-keff}
\boxed{
K_0^{\rm eff}=K_0-C\Omega_0^{-2}C^T,
\qquad
M_0^{\rm eff}=M_0+C\Omega_0^{-4}C^T,
\qquad
N_0^{\rm eff}=C\Omega_0^{-6}C^T}.
\end{equation}
Stable matter support therefore lowers static stiffness, increases effective inertia, and generates the next even moment.

\paragraph{2. Grouped real \(P_2\) reduction.}
For grouped lane \(A\in\{20,21,22\}\), let \(q_A\) be the wall amplitude and \(X_{A\alpha}\) the stable BdG support coordinates.  The reduced Lagrangian is
\begin{equation}
\label{eq:app-stage003-p2-lagrangian}
L_{\rm red}^{(2)}=
\sum_A\left[
\frac12M_A\dot q_A^2-\frac12K_Aq_A^2
+\frac12\sum_\alpha(\dot X_{A\alpha}^2-\varpi_{A\alpha}^2X_{A\alpha}^2)
+\sum_\alpha g_{A\alpha}q_AX_{A\alpha}
\right].
\end{equation}
Eliminating the support coordinates gives
\begin{equation}
\label{eq:app-stage003-p2-kernel}
\boxed{
D_A^{\rm eff}(\omega)
=K_A-M_A\omega^2-
\sum_\alpha\frac{g_{A\alpha}^2}{\varpi_{A\alpha}^2-\omega^2},
\qquad
A\in\{20,21,22\}}.
\end{equation}
At low frequency,
\begin{equation}
D_A^{\rm eff}(\omega)=d_{0A}^{\rm eff}+d_{2A}^{\rm eff}\omega^2+d_{4A}^{\rm eff}\omega^4+O(\omega^6),
\end{equation}
where
\begin{align}
\label{eq:app-stage003-d0d2d4}
d_{0A}^{\rm eff}&=K_A-\sum_\alpha\frac{g_{A\alpha}^2}{\varpi_{A\alpha}^2},\\
d_{2A}^{\rm eff}&=-\left(M_A+\sum_\alpha\frac{g_{A\alpha}^2}{\varpi_{A\alpha}^4}\right),\\
d_{4A}^{\rm eff}&=-\sum_\alpha\frac{g_{A\alpha}^2}{\varpi_{A\alpha}^6}.
\end{align}
The executable audit closes finite-mode witness cases explicitly; the displayed
\(\sum_\alpha\) formulas then follow by linear superposition because each stable
BdG mode contributes additively to the Schur complement.

\paragraph{3. One-wall-mode pole shift.}
For the two-coordinate prototype
\begin{equation}
L=\frac12M\dot q^2-\frac12Kq^2+\frac12\dot X^2-\frac12\varpi^2X^2+gqX,
\end{equation}
the dispersion relation is
\begin{equation}
\label{eq:app-stage003-dispersion}
(K-M\omega^2)(\varpi^2-\omega^2)-g^2=0.
\end{equation}
Writing \(\Omega_\eta^2=K/M\), the exact poles are
\begin{equation}
\label{eq:app-stage003-poles}
\boxed{
\omega_\pm^2=\frac{\Omega_\eta^2+\varpi^2
\pm\sqrt{(\Omega_\eta^2-\varpi^2)^2+4g^2/M}}{2}}.
\end{equation}
For \(\varpi^2>\Omega_\eta^2\) and weak coupling,
\begin{equation}
\label{eq:app-stage003-pole-shift}
\boxed{
\delta\Omega_\eta^2=-\frac{g^2/M}{\varpi^2-\Omega_\eta^2}+O(g^4),
\qquad
\delta\varpi^2=+\frac{g^2/M}{\varpi^2-\Omega_\eta^2}+O(g^4)}.
\end{equation}
This is the first explicit reduced formula for how matter support shifts a wall pole.

\paragraph{4. Grouped isotropy diagnostic.}
For any grouped triple \(x_A=(x_{20},x_{21},x_{22})\), define
\begin{equation}
\label{eq:app-stage003-trace-anomaly}
\bar x=\frac{x_{20}+2x_{21}+2x_{22}}{5},
\qquad
a_x=\frac{2x_{20}-x_{21}-x_{22}}{10},
\qquad
b_x=\frac{x_{21}-x_{22}}{2}.
\end{equation}
Applying this to \(d_{2A}^{\rm eff}\) gives grouped conservative anisotropy defects.  If
\begin{equation}
K_{20}=K_{21}=K_{22},
\quad
M_{20}=M_{21}=M_{22},
\quad
g_{20,\alpha}=g_{21,\alpha}=g_{22,\alpha},
\quad
\varpi_{20,\alpha}=\varpi_{21,\alpha}=\varpi_{22,\alpha},
\end{equation}
then all grouped \(d_{nA}^{\rm eff}\) are lane-independent and
\begin{equation}
\boxed{a_x=b_x=0}.
\end{equation}
Conversely, a nonzero defect must come from anisotropy in wall data, matter support frequencies, or wall--matter overlaps.

\stagefield{Output}{Stage~003 outputs the scalar Schur complement \eqref{eq:app-stage003-scalar-kernel}, the grouped kernels \eqref{eq:app-stage003-p2-kernel}, the low-frequency coefficients \eqref{eq:app-stage003-d0d2d4}, the exact pole shift \eqref{eq:app-stage003-pole-shift}, and the first microscopic grouped-isotropy criterion.}

\stagefield{Checks}{\begin{verificationchecklist}
\item Expanding \((\Omega^2-\omega^2)^{-1}=\Omega^{-2}+\omega^2\Omega^{-4}+\omega^4\Omega^{-6}+\cdots\) gives \eqref{eq:app-stage003-keff} and \eqref{eq:app-stage003-d0d2d4}.
\item The sign of the static correction is softening: the Schur complement subtracts \(g^2/\varpi^2\).
\item The pole shift has the expected avoided-crossing sign: the lower wall-like pole moves down when coupled to a higher stable matter mode.
\item The isotropy check follows by lane equality of all microscopic inputs.
\end{verificationchecklist}}

\stagefield{Downstream use}{Stage~004 adds Maxwell/mixed self-energies to the same operator structure.  Stage~005 converts the \(D_A\) operator moments into normalized grouped response moments.  Stage~006 collects the Stage~003 BdG moments as \(B_{A0},B_{A2},B_{A4}\).}

\clearpage
\section[Stage 004]{Stage 004: Localized Maxwell/Mixed Response and the First Outgoing Quadrupole Bridge}
\label{stage:004}

\stagefield{Part / anchor}{Part I; primary anchors \resultanchor{MTDC-T3} and \resultanchor{MTDC-T4}.}
\claimstatus{\StatusExactClosure{} inside the one-lane reduced Maxwell/mixed model.  The model is \StatusReduced{} relative to the full localized Maxwell PDE.  Actual branch realization of the transfer coefficient is \StatusOpen{}.}

\stagefield{Purpose}{Stage~004 introduces the first sector capable of carrying a passive/outgoing channel: the localized Maxwell/mixed block.  It shows that mixed fields are gauge-invariant, derives the conservative Maxwell/mixed self-energy, attaches an outgoing port to the mixed coordinate, and obtains the first wall-level \(i\omega^5\) quadrupole coefficient.}

\stagefield{Inputs}{The stage imports the exact localized Maxwell equation, the mixed gauge invariants \(E_w=F_{w0}\) and \(C_a=F_{aw}\), the Stage~003 wall operator, and the retarded-branch result that the surviving universal point-particle route is the outgoing \(l=2\) quadrupole branch.}

\stagefield{Verification}{SymPy audit: \StageFile{scripts/moving_throat_pde_stage004_maxwell_mixed_sympy_audit.py}; transcript: \StageFile{scripts/output/moving_throat_pde_stage004_maxwell_mixed_sympy_audit.txt}.  Mathematica audit: \StageFile{mathematica/moving_throat_pde_stage004_maxwell_mixed_mathematica_audit.wl}; transcript: \StageFile{mathematica/output/moving_throat_pde_stage004_maxwell_mixed_mathematica_audit.txt}.  Shared numerical stress test with Stage~003: \StageFile{scripts/numerical/stage003_004_foundational_stress.py}, \StageFile{scripts/numerical/output/stage003_004_foundational_stress.txt}, \StageFile{mathematica/numerical/stage003_004_foundational_stress.wl}, and \StageFile{mathematica/numerical/output/stage003_004_foundational_stress.txt}.}

\subsection*{Derivation}

\paragraph{1. Mixed gauge invariants.}
Under
\begin{equation}
A_0\mapsto A_0-\partial_t\chi,
\qquad
A_a\mapsto A_a+\partial_a\chi,
\qquad
A_w\mapsto A_w+\partial_w\chi,
\end{equation}
the mixed fields
\begin{equation}
\label{eq:app-stage004-mixed-fields}
E_w=-\partial_tA_w-\partial_wA_0,
\qquad
C_a=\partial_aA_w-\partial_wA_a
\end{equation}
are gauge invariant.  Therefore a reduced coordinate representing the mixed \(A_w/F_{\mu w}/J^w\) sector is not a gauge artifact.

\paragraph{2. One-lane Maxwell/mixed reduced Lagrangian.}
For one harmonic lane \(l\), let \(Q_l\) be the wall amplitude, \(A_l\) a brane-like localized Maxwell coordinate, and \(W_l\) a mixed-sector coordinate.  The reduced Lagrangian is
\begin{align}
\label{eq:app-stage004-reduced-lagrangian}
L_l={}&\frac12M_l\dot Q_l^{\,2}-\frac12K_lQ_l^2
+\frac12\dot A_l^{\,2}-\frac12\Omega_{A,l}^2A_l^2
+\frac12\dot W_l^{\,2}-\frac12\Omega_{W,l}^2W_l^2 \\
&+R_lA_lW_l+g_{A,l}Q_lA_l+g_{W,l}Q_lW_l.
\end{align}
Use the frequency convention \(e^{-i\omega t}\) and define
\begin{equation}
A_l(\omega)=\Omega_{A,l}^2-\omega^2,
\qquad
W_l(\omega)=\Omega_{W,l}^2-\omega^2,
\qquad
\Delta_l(\omega)=A_l(\omega)W_l(\omega)-R_l^2.
\end{equation}
Eliminating \((A_l,W_l)\) gives the conservative self-energy
\begin{equation}
\label{eq:app-stage004-self-energy}
\boxed{
\Sigma_l^{({\rm EM}+{\rm mix},{\rm cons})}(\omega)
=\frac{g_{A,l}^2W_l(\omega)+2g_{A,l}g_{W,l}R_l+g_{W,l}^2A_l(\omega)}{\Delta_l(\omega)}}.
\end{equation}
The wall operator is
\begin{equation}
\label{eq:app-stage004-wall-operator-cons}
D_l^{\rm cons}(\omega)=K_l-M_l\omega^2-\Sigma_l^{({\rm EM}+{\rm mix},{\rm cons})}(\omega).
\end{equation}

\paragraph{3. Low-frequency conservative coefficients.}
At \(\omega=0\), define
\begin{equation}
\Delta_0=\Omega_A^2\Omega_W^2-R^2,
\qquad
S_2=\Omega_A^2+\Omega_W^2,
\end{equation}
\begin{equation}
N_0^{\rm cons}=g_A^2\Omega_W^2+2g_Ag_WR+g_W^2\Omega_A^2,
\qquad
G_2=g_A^2+g_W^2.
\end{equation}
Then
\begin{equation}
\Sigma^{\rm cons}(\omega)=z_0+z_2\omega^2+z_4\omega^4+O(\omega^6),
\end{equation}
with
\begin{align}
\label{eq:app-stage004-z-coeffs}
z_0&=\frac{N_0^{\rm cons}}{\Delta_0},\\
z_2&=\frac{N_0^{\rm cons}S_2-G_2\Delta_0}{\Delta_0^2},\\
z_4&=\frac{N_0^{\rm cons}(S_2^2-\Delta_0)-S_2G_2\Delta_0}{\Delta_0^3}.
\end{align}
These coefficients later become the conservative Maxwell/mixed entries \(Z_{A0},Z_{A2},Z_{A4}\).

\paragraph{4. Outgoing-port dressing and transfer factor.}
Attach the outgoing port to the mixed coordinate by replacing
\begin{equation}
W_l(\omega)\longmapsto W_l(\omega)-\Pi_l^{\rm out}(\omega).
\end{equation}
To first order in \(\Pi_l^{\rm out}\),
\begin{equation}
\Sigma_l^{\rm full}(\omega)
=\Sigma_l^{\rm cons}(\omega)+\Pi_l^{\rm out}(\omega)N_l(\omega)+O\bigl((\Pi_l^{\rm out})^2\bigr),
\end{equation}
where the exact transfer factor is
\begin{equation}
\label{eq:app-stage004-transfer-factor}
\boxed{
N_l(\omega)=
\frac{\bigl[A_l(\omega)g_{W,l}+R_lg_{A,l}\bigr]^2}{\bigl[A_l(\omega)W_l(\omega)-R_l^2\bigr]^2}}.
\end{equation}
At zero frequency,
\begin{equation}
\label{eq:app-stage004-n0-positive}
\boxed{
N_l(0)=\frac{(\Omega_{A,l}^2g_{W,l}+R_lg_{A,l})^2}{(\Omega_{A,l}^2\Omega_{W,l}^2-R_l^2)^2}\ge0}.
\end{equation}
The outgoing sign is therefore not inserted by hand.  The reduced internal block transfers the outgoing port with a nonnegative conservative factor.

\paragraph{5. Compact outgoing \(l=2\) fingerprint.}
The normalized compact outgoing \(l=2\) response has the low-frequency form
\begin{equation}
\label{eq:app-stage004-outgoing-fingerprint}
\widehat Y_2^{\rm out}(\omega)
=1+\frac{a^2\omega^2}{9c_s^2}
+\frac{4a^4\omega^4}{81c_s^4}
+i\frac{a^5\omega^5}{27c_s^5}
+O(\omega^6).
\end{equation}
Thus
\begin{equation}
\label{eq:app-stage004-gamma-port}
\Gamma_2^{\rm port}=\frac{a^5}{27c_s^5}.
\end{equation}
In the wall-operator convention \(D=K-M\omega^2-\Sigma\), a positive-imaginary outgoing port contributes
\begin{equation}
\label{eq:app-stage004-wall-odd}
\boxed{
\delta D_2^{\rm odd}(\omega)
=-iN_2(0)\frac{a^5}{27c_s^5}\omega^5+O(\omega^7)}.
\end{equation}
The same branch appears with the corresponding positive admittance sign when the response is normalized.

\paragraph{6. Scalar compatibility check.}
If the scalar/breathing lane enters the port-active mixed sector only through a derivative coupling
\begin{equation}
g_{W,0}(\omega)=\eta\omega,
\end{equation}
with no non-derivative direct coupling into the same port-active combination, then
\begin{equation}
N_0(\omega)=O(\omega^2).
\end{equation}
A scalar outgoing law \(\Pi_0^{\rm out}(\omega)=i\gamma_1\omega+\cdots\) then induces only
\begin{equation}
\delta D_0^{\rm odd}(\omega)=O(i\omega^3),
\end{equation}
so the dangerous \(i\omega\) scalar wall-level term is delayed.  This is a compatibility mechanism, not yet a full scalar no-go theorem.

\stagefield{Output}{Stage~004 outputs the conservative self-energy \eqref{eq:app-stage004-self-energy}, its low-frequency coefficients \eqref{eq:app-stage004-z-coeffs}, the outgoing transfer factor \eqref{eq:app-stage004-transfer-factor}, the nonnegative static transfer \eqref{eq:app-stage004-n0-positive}, and the wall-level \(i\omega^5\) quadrupole coefficient \eqref{eq:app-stage004-wall-odd}.}

\stagefield{Checks}{\begin{verificationchecklist}
\item Gauge invariance of \(E_w\) and \(C_a\) follows by direct substitution of the gauge transformation.
\item The Schur complement of the \((A,W)\) block gives \eqref{eq:app-stage004-self-energy}.
\item The outgoing transfer coefficient is a square over a squared determinant at \(\omega=0\), hence nonnegative.
\item The compact outgoing fingerprint has the required \(l=2\) odd order \(\omega^{2l+1}=\omega^5\).
\item Derivative scalar coupling delays a scalar \(i\omega\) port to a wall-level \(i\omega^3\) contribution.
\end{verificationchecklist}}

\stagefield{Downstream use}{Stage~005 uses \(N_A(\omega)\) as the outgoing-transfer factor in each grouped lane.  Stage~006 collects the conservative coefficients from \eqref{eq:app-stage004-z-coeffs} into \(Z_{A0},Z_{A2},Z_{A4}\) and the transfer moments into \(N_{A0},N_{A2},N_{A4}\).  Later retarded-closure stages use \eqref{eq:app-stage004-outgoing-fingerprint} and \eqref{eq:app-stage004-wall-odd}.}

\clearpage
\section[Stage 005]{Stage 005: Covariant projected Maxwell law}
\label{stage:005}

\stagefield{Part / anchor}{Part I; primary anchor \resultanchor{MTDC-T4}.}
\claimstatus{This stage is \StatusExact{} for the displayed algebraic identities inside the declared parent/projection data, with \StatusReduced{} or \StatusExactClosure{} inherited where the local text invokes the matched reduction, mouth-kernel closure, parent-action packet, or one-port normal form.}

\paragraph{Source and audit.}
This stage corresponds to
\StageFile{scripts/moving_throat_pde_stage005_projected_maxwell_covariant_sympy_audit.py}.
It is the covariant projection of the localized parent Maxwell equation.

\paragraph{Parent law.}
The localized parent equation is
\begin{equation}
\label{eq:stage005-parent-maxwell}
\partial_N\!\left(Z(w)F^{NM}\right)=\mu_0J^M .
\end{equation}
For brane-facing components \(M=\mu\), apply a normalized observation kernel
\(W(w)\) and integrate in \(w\).  The result is
\begin{equation}
\label{eq:stage005-projected-maxwell}
\partial_\nu\!\int WZ F^{\nu\mu}\,dw
=
\mu_0\int WJ^\mu\,dw
-\int W\partial_w(ZF^{w\mu})\,dw.
\end{equation}
Using \eqref{eq:stage004-projection-ibp},
\begin{equation}
\label{eq:stage005-projected-maxwell-expanded}
\partial_\nu\!\int WZ F^{\nu\mu}\,dw
=
\mu_0\int WJ^\mu\,dw
-[WZF^{w\mu}]_{\partial}
+\int (\partial_wW)ZF^{w\mu}\,dw .
\end{equation}

\paragraph{Projected continuity.}
The same projection applied to parent current conservation gives
\begin{equation}
\label{eq:stage005-projected-continuity}
\partial_\mu\!\int WJ^\mu\,dw
=-[WJ^w]_{\partial}+\int(\partial_wW)J^w\,dw .
\end{equation}
Thus leakage is not a notation choice.  It is the exact open-system balance
term induced by observing a localized parent current through \(W\).

\paragraph{Output.}
Stage~005 exports the projected inhomogeneous law
\eqref{eq:stage005-projected-maxwell}--\eqref{eq:stage005-projected-maxwell-expanded}
and projected continuity \eqref{eq:stage005-projected-continuity}.

\clearpage
\section[Stage 006]{Stage 006: Vector projected Maxwell split}
\label{stage:006}

\stagefield{Part / anchor}{Part I; primary anchor \resultanchor{MTDC-T4}.}
\claimstatus{This stage is \StatusExact{} for the displayed algebraic identities inside the declared parent/projection data, with \StatusReduced{} or \StatusExactClosure{} inherited where the local text invokes the matched reduction, mouth-kernel closure, parent-action packet, or one-port normal form.}

\paragraph{Source and audit.}
This stage corresponds to
\StageFile{scripts/moving_throat_pde_stage006_projected_maxwell_vector_sympy_audit.py}.
It checks the \(3+1\) vector signs, the measured-field/source-flux split, and
the boundary discharge terms.

\paragraph{Measured and flux fields.}
The homogeneous equations use the measured projected fields
\begin{equation}
\label{eq:stage006-measured-fields}
\mathbf E_{\rm meas}=\int W\mathbf E\,dw,
\qquad
\mathbf B_{\rm meas}=\int W\mathbf B\,dw,
\end{equation}
whereas the inhomogeneous equations use the source-coupled flux fields
\begin{equation}
\label{eq:stage006-flux-fields}
\mathbf D_{\rm flux}=\int WZ\mathbf E\,dw,
\qquad
\mathbf H_{\rm flux}=\int WZ\mathbf B\,dw.
\end{equation}
Identifying these two pairs is an additional closure assumption.  It is not
part of the parent projection.

\paragraph{Projected vector laws.}
With the convention \(F^{i0}=E_i\), the projected Faraday and Ampere signs are
fixed by the audit:
\begin{equation}
\partial_t\mathbf B_{\rm meas}+\nabla\times\mathbf E_{\rm meas}=0,
\end{equation}
\begin{equation}
\label{eq:stage006-ampere}
-\nabla\times\mathbf H_{\rm flux}-\partial_t\mathbf D_{\rm flux}
+\mathbf L_{\rm mix}=\mu_0\mathbf J_{\rm proj}.
\end{equation}
Here \(\mathbf L_{\rm mix}\) is the vector form of the transverse leakage term
in \eqref{eq:stage005-projected-maxwell-expanded}.

\paragraph{Output.}
Stage~006 exports the field split
\eqref{eq:stage006-measured-fields}--\eqref{eq:stage006-flux-fields}
and the vector-law sign convention used by the rest of the projected Maxwell
block.

\endgroup

\section{Part I appendix closure and handoff to Part II}
\label{sec:app-part01-closure}

Stages~001--006 complete the first verification block.  They do not solve the nonlinear moving-throat PDE, but they remove the main algebraic ambiguity that existed before the ledger began.  The reduced response variables used later are not arbitrary symbols: they are the low-frequency coefficients of a coupled wall/BdG/Maxwell/mixed Schur-complement system, organized by exact grouped projectors and normalized by a definite outgoing \(l=2\) fingerprint.

The handoff to Part~II is therefore precise.  What remains to be checked next is not the abstract existence of grouped \(P_2\) lanes, nor whether the electromagnetic sector should be projected before it is reduced.  Those now exist inside the reduced moving-throat construction.  The next question is whether explicit radial/axial overlap integrals on an isotropic branch force the three grouped lanes to agree, and how the first controlled anisotropy enters when the branch is weakly axisymmetric.  That is the job of Stages~007--019.
